\documentclass[12pt]{article}
\usepackage{sbc-template}
\usepackage{graphicx,url}
\usepackage[utf8]{inputenc}
\usepackage[brazilian]{babel}
\usepackage{xcolor}
\usepackage{courier}

% Figuras ficam em ../imagens (raiz do projeto).
\graphicspath{{../imagens/}}

\sloppy
\title{Reconhecimento Facial em Conjunto Aberto para Controle de Acesso\\com Poucas Amostras de Cadastro}
\author{Gustavo Cesar Bento Laurindo\inst{1}, Gustavo de Oliveira Gimenes\inst{1},\\ Lucas Vasconcelos Fujii\inst{1}, Mairon Cesar de Paula Faria\inst{1}}
\address{Departamento de Computa\c{c}\~ao -- Universidade Federal de S\~ao Carlos (UFSCar)\\
S\~ao Carlos -- SP -- Brasil
\email{\{gustavocbl, ggimenes, lucasfujii, maironfaria\}@estudante.ufscar.br}}
\date{}
\begin{document}
\maketitle

\begin{abstract}
TODO.
\end{abstract}

\begin{resumo}
TODO.
\end{resumo}

\section{Introdução}
% Controle de acesso facial; por que classificação fechada falha; conjunto aberto.

\section{Trabalhos relacionados}
% FaceNet, ArcFace, SFace; reconhecimento em conjunto aberto; one-class.
TODO: FaceNet~\cite{schroff2015facenet}, ArcFace~\cite{deng2019arcface}, LFW~\cite{huang2007lfw}.

\section{Dados}
% Rostos do grupo (sessões distintas, consentimento, LGPD) e desconhecidos (LFW/CelebA).
% Divisão por sessão (integrantes) e por identidade (desconhecidos).

\section{Método}
\subsection{Pipeline: detecção, alinhamento e embedding}
\subsection{CNNs avaliadas}
% Do zero (ResNet18/CelebA), pré-treinadas, ajuste fino com perda ArcFace.
\subsection{Regras de decisão em conjunto aberto}
% Limiar no centróide, kNN, One-Class SVM, SVM com negativos, softmax com limiar.
\subsection{Calibração do limiar e métricas}
% FAR, FRR, EER, curva DET, identificação a FAR fixo.

\section{Resultados}
\subsection{Comparação entre CNNs e regras de decisão}
\subsection{Número de amostras de cadastro}
\subsection{Condições adversas e impostores difíceis}
\subsection{Cadastro posterior e crescimento da galeria}

\section{Discussão e limitações}
% Ataques de apresentação fora do escopo; viés; privacidade.

\section{Conclusão}

\bibliographystyle{sbc}
\bibliography{referencias}
\end{document}
